% \iffalse meta-comment
%
% hit-linebreak.dtx 是断行层：LuaTeX 上照 Word 的贪心断行
%
% \fi
%
% \section{断行层}
%
% \changes{v3.2a}{2026/09/15}{新增断行层：\hologo{LuaLaTeX} 上不再用 TeX 的最优断行，
%   照 Word 的贪心算法逐行首次适配、装不下就压标点，规则全是量出来的}
% \changes{v3.2a}{2026/09/15}{兼容模式 11、12、14、15 各量一遍：学校的 doc 原件都是 11，
%   默认改成 11；\option{word-compat} 收四个值}
% Word 排一段是逐行贪心：这一行能装几个装几个，下一个字装不下，先看能不能把
% 行内的标点压一压塞进来，塞不进就换行；行首不能是闭号，行尾不能是开号，碰上了
% 先压标点再悬挂。TeX 是全段最优，两边永远对不上。\hologo{LuaLaTeX} 的
% \texttt{linebreak\_filter} 回调把整段的节点表交给 Lua，Lua 自己切行、自己定
% 每处胶与 kern 的宽、自己装盒，TeX 只管把行盒往下摞。
%
% 规则不是猜的：拿 \texttt{python-docx} 自造几百个探针段落让 Word 排，读导出的
% PDF 逐字量，再拿学校的本科范例逐行校对。压缩总额 $T = p\,K/(K+1)$（$p$ 是
% 格宽、$K$ 是行内压缩位的权重和）这条对 $K = 1$ 到 $10$ 全对，误差在
% $0.02$\,bp 之内。全部规则写在 Lua 文件头的注释里，量法与探针记在维护者的
% 笔记里。
%
% \subsection{开关}
%
% \option{layout}/\option{linebreak} 取 \texttt{word}（默认）或 \texttt{tex}。
% \texttt{word} 只在 \hologo{LuaLaTeX} 且走 Word 版面时生效，
% \hologo{XeLaTeX} 上没有这条路，写了也当 \texttt{tex}。
%
% \option{layout}/\option{word-compat} 是 Word 文件的兼容模式，就是 docx 的
% \file{settings.xml} 里 \texttt{compatibilityMode} 那个数：\texttt{11}（Word 2003）、
% \texttt{12}（2007）、\texttt{14}（2010）、\texttt{15}（2013 起）。学校发的原件几乎都是
% \texttt{doc}，用新版 Word 原样另存都是 11 且带 2003 的三个开关；学校自己发的少数
% docx/dotx 是 12 或 14；15 只出现在用户点过「转换」的文件里。所以默认 \texttt{11}。
%
% 四档里 11、12、14 的断行规则一样，只差行末闭号一处没有行内位时挂半格（11）还是
% 整格（12、14）；15 另一套，能多塞汉字。规则见 Lua 文件头。
%    \begin{macrocode}
%<*shared-linebreak>
\bool_new:N \g__hit_linebreak_word_bool
\bool_gset_true:N \g__hit_linebreak_word_bool
\int_new:N \g__hit_linebreak_compat_int
\int_gset:Nn \g__hit_linebreak_compat_int { 11 }
\keys_define:nn { hit / layout }
  {
    linebreak .choice: ,
    linebreak / word .code:n = { \bool_gset_true:N  \g__hit_linebreak_word_bool } ,
    linebreak / tex  .code:n = { \bool_gset_false:N \g__hit_linebreak_word_bool } ,
    linebreak .value_required:n = true ,
    word-compat .choice: ,
    word-compat / 11 .code:n = { \int_gset:Nn \g__hit_linebreak_compat_int { 11 } } ,
    word-compat / 12 .code:n = { \int_gset:Nn \g__hit_linebreak_compat_int { 12 } } ,
    word-compat / 14 .code:n = { \int_gset:Nn \g__hit_linebreak_compat_int { 14 } } ,
    word-compat / 15 .code:n = { \int_gset:Nn \g__hit_linebreak_compat_int { 15 } } ,
    word-compat .value_required:n = true ,
  }
%    \end{macrocode}
%
% \subsection{装回调}
%
% 属性 \cs{hit@linebreak@attr} 是逐段的门：段首节点上它为 0 的段交回 TeX 断行。
% 正文默认为 1；不想让 Lua 碰的环境把它清零。Lua 自己另有几道门：有 \cs{parshape}
% 或悬挂缩进的、\cs{leftskip} 非零或 \cs{rightskip} 可伸的（居中、靠左）、
% \cs{hsize} 不等于版心宽的（表格、小页），都交回 TeX。
%
% 回调在 \texttt{begindocument/end} 装：Lua 模块要用 \pkg{luatexja} 的表，
% 那时它肯定装好了；格宽与正文字号从 \cs{g\_hit\_msword\_char\_dim} 与
% \cs{g\_hit\_body\_size\_dim} 两个寄存器取，也都定下来了。
% Lua 文件装在 \file{tex/latex/hithesis/} 下，跟类放一起；\texttt{require}
% 只找 \texttt{LUAINPUTS}（\file{scripts/} 与 \file{tex/luatex/}），所以按
% \texttt{tex} 格式用 \texttt{kpse} 找到再 \texttt{dofile}。
%    \begin{macrocode}
\sys_if_engine_luatex:T
  {
    \newattribute \hit@linebreak@attr
    \hit@linebreak@attr = 1
    \AddToHook { begindocument / end }
      {
        \bool_if:NT \g__hit_linebreak_word_bool
          {
            \__hit_layout_msword_active:T
              {
                \tex_directlua:D
                  {
                    hithesis = hithesis~or~{}
                    hithesis.linebreak = dofile(
                      kpse.find_file("hit-linebreak.lua", "tex")~
                      or~kpse.find_file("hit-linebreak.lua", "lua"))
                    hithesis.linebreak.compat = \int_use:N \g__hit_linebreak_compat_int
                    hithesis.linebreak.install(luatexbase.attributes["hit@linebreak@attr"])
                  }
              }
          }
      }
  }
%</shared-linebreak>
%    \end{macrocode}
%
% \subsection{Lua 模块}
%
% 下面整段解出来是 \file{hit-linebreak.lua}。docstrip 逐行搬运，\texttt{\%} 开头的
% 行当注释丢掉，其余原样，所以 Lua 代码里的行首不能有 \texttt{\%}。
%    \begin{macrocode}
%<*lua>
-- hit-linebreak.lua：hithesis 的断行层，照 Word 的贪心断行（LuaTeX linebreak_filter）
-- 规则全是自造 docx 让 Word 排、逐行量出来的：
--   自然宽：汉字/开号一格 p；闭号一格，后接标点则半格；西文字形宽乘 p/em（Word 的
--     「中西文字宽平衡」）；汉字与西文字母数字之间自动间距 3.25bp；西文接闭号 0.27bp；
--     两边都是中文的空格占半格且不伸缩，其余空格 3.5bp；—～· 这类全角符号当西文串里的字，不加间距。
--   塞汉字：压缩位 K = 行内闭号（后接闭号的算 0.2）+ 开号（前面是汉字 1、是标点 0.2）。
--     模式 15：总额 T = p·K/(K+1)，各位均摊，行里有西文块再按 K 扣一点（量出来的表）；
--     模式 11/12/14：有位就给半格，没有不给，与 K 无关，西文不扣。
--     行末是开号：再加开号自己的左半空白，装不进就把开号拉到下一行。
--   塞闭号（禁则）：模式 15 一处行内位都没有时整个闭号挂出去，有 K 个则 p + p·(K+1)/(K+3)；
--     模式 11（带 doNotWrapTextWithPunct）挂半格、12/14 挂整格，行内位各半格、至多两个。
--   松行：余量先给汉西隙（每处至多 2.8bp）与西文空格，再均摊到汉字缝；贴格段伸到整数格。
--   兼容模式的来历：学校的 doc 原件用新版 Word 另存都是 11 且带 2003 那三个开关；
--   学校自己发的 docx/dotx 是 12 或 14；15 只出现在用户「转换」过的文件里。
local M = {}
M.version = "v3.2a"
local node, tex = node, tex
local traverse = node.traverse
local GLYPH, GLUE, KERN, PENALTY, RULE, HLIST, VLIST, WHATSIT, MATH, DISC, LOCALPAR =
  node.id("glyph"), node.id("glue"), node.id("kern"), node.id("penalty"), node.id("rule"),
  node.id("hlist"), node.id("vlist"), node.id("whatsit"), node.id("math"), node.id("disc"), node.id("local_par")
local INS, MARK, ADJUST = node.id("ins"), node.id("mark"), node.id("adjust")
local attr_icflag = luatexbase.attributes['ltj@icflag']
local IC = luatexja.icflag_table

local function set_of(s)
  local t = {}
  for _, c in utf8.codes(s) do t[c] = true end
  return t
end
local closers = set_of("。，、；：？！）》〉”’」』】］．")
local openers = set_of("（《〈「『【“‘［")
-- 西文标点的禁则：行首不能是这些，行尾不能是那些
local ascii_close = set_of(".,;:!?)]}%")
local ascii_open = set_of("([{")
-- 全角符号里当西文处理的：波浪线、破折号、间隔号、连接号
local latinish = set_of("～—–·")

M.trace = false
M.compat = 11   -- Word 兼容模式：11/12/14 一档（学校发的文件），15 另一档（转换过的 docx）
M.snap = false  -- 段落贴行网格（Word 的 snapToGrid）
M.ari0 = true   -- 模式 ≤14：adjustRightInd=0 时装到版心，否则整数格
local function trace(...)
  if M.trace then texio.write_nl("term and log", "HG| " .. string.format(...)) end
end

local function is_jfont(f)
  return luatexja.jfont.font_metric_table[f] ~= nil
end
local function icflag(n)
  return node.get_attribute(n, attr_icflag) or 0
end
local function is_jfmglue(n)
  local ic = icflag(n)
  return ic >= IC.FROM_JFM and ic < IC.KANJI_SKIP
end
local function is_kanjiskip(n)
  local ic = icflag(n)
  return ic == IC.KANJI_SKIP or ic == IC.KANJI_SKIP_JFM
end
local function is_xkanjiskip(n)
  local ic = icflag(n)
  return ic == IC.XKANJI_SKIP or ic == IC.XKANJI_SKIP_JFM
end
local function isdigit(c)
  return c and c >= 48 and c <= 57
end
local function is_alnum(c)
  if not c then return false end
  return (c >= 48 and c <= 57) or (c >= 65 and c <= 90) or (c >= 97 and c <= 122)
    or (c >= 0xC0 and c <= 0x24F) or (c >= 0x370 and c <= 0x3FF)
end

-- 把节点表切成条目：0 汉字 1 闭号 2 开号 3 西文串 4 空格
-- 西文串记字形宽之和 gw、字数 gn、首末字符；串里的字距 kern 不计宽（渲染时重设）
local function tokenize(head)
  local items = {}
  local run = nil
  local function flush_run()
    if run then items[#items + 1] = run; run = nil end
  end
  local function put(kind, first, last, w, c)
    flush_run()
    items[#items + 1] = { kind = kind, first = first, last = last, w = w, c = c }
    return items[#items]
  end
  local function attach(n)
    if run then run.last = n elseif #items > 0 then items[#items].last = n end
  end
  local function addglyph(n, c, jglyph)
    if not run then run = { kind = 3, first = n, last = n, gw = 0, gn = 0, head = c, tail = c, kerns = {}, glyphs = {} } end
    run.last = n; run.gw = run.gw + n.width; run.gn = run.gn + 1; run.tail = c
    run.glyphs[#run.glyphs + 1] = n
    if jglyph then run.hasj = true end
  end
  local n = head
  local indent = 0
  local strut_h, strut_d = 0, 0
  local lead = nil   -- 第一个条目之前的节点（链接起点这类 whatsit），要归进第一行
  while n do
    local id = n.id
    local nxt = n.next
    if #items == 0 and not run and not lead and id ~= LOCALPAR and not (id == HLIST and n.subtype == 3) then lead = n end
    if id == LOCALPAR then
    elseif id == HLIST and n.subtype == 3 then
      -- 缩进盒不进任何条目（渲染时另放一个），它前面的东西也不要
      indent = n.width
      lead = nil
    elseif id == GLYPH then
      local c = n.char
      if is_jfont(n.font) and not latinish[c] then
        if closers[c] then put(1, n, n, n.width, c)
        elseif openers[c] then put(2, n, n, n.width, c)
        else put(0, n, n, n.width, c) end
      else
        -- 连字符后面可断（数字-数字不断）；波浪线、破折号后面可断
        addglyph(n, c, is_jfont(n.font))
        if c == 45 or latinish[c] then flush_run() end
      end
    elseif id == HLIST and icflag(n) == IC.PACKED then
      local g = n.head
      while g and g.id ~= GLYPH do g = g.next end
      local c = g and g.char or 0
      if closers[c] then put(1, n, n, n.width, c)
      elseif openers[c] then put(2, n, n, n.width, c)
      else put(0, n, n, n.width, c) end
    elseif id == GLUE then
      if is_kanjiskip(n) or is_xkanjiskip(n) then
        attach(n)
      elseif is_jfmglue(n) then
        if #items > 0 and not run and items[#items].kind == 1 and not items[#items].jfmglue then
          items[#items].last = n; items[#items].jfmglue = n
        else
          attach(n)
        end
      elseif n.subtype == 13 or n.subtype == 12 or n.subtype == 0 then
        put(4, n, n, n.width, nil)
      else
        attach(n)
      end
    elseif id == KERN then
      if run then run.last = n; run.kerns[#run.kerns + 1] = n else attach(n) end
    elseif id == DISC then
      -- 明写的连字符在 LuaTeX 里是 disc（replace 里放着字形）：计宽，后面可断
      local w, c = 0, nil
      for g in traverse(n.replace) do
        if g.id == GLYPH then w = w + g.width; c = g.char elseif g.id == KERN then w = w + g.kern end
      end
      if c then
        if not run then run = { kind = 3, first = n, last = n, gw = 0, gn = 0, head = c, tail = c, kerns = {}, glyphs = {} } end
        run.last = n; run.gw = run.gw + w; run.gn = run.gn + 1; run.tail = c
        run.discs = run.discs or {}; run.discs[#run.discs + 1] = n
        if c == 45 then flush_run() end
      else
        attach(n)
      end
    elseif id == PENALTY then
      attach(n)
    elseif id == RULE then
      if n.width == 0 then
        if n.height > strut_h then strut_h = n.height; strut_d = n.depth end
        attach(n)
      else
        put(3, n, n, n.width, nil).rigid = true
      end
    elseif id == MATH then
      local last = n; local w = 0
      local m = n.next
      while m and m.id ~= MATH do w = w + (m.width or 0); last = m; m = m.next end
      if m then last = m end
      flush_run()
      items[#items + 1] = { kind = 3, first = n, last = last, w = w, math = true, rigid = true }
      nxt = last.next
    elseif id == HLIST or id == VLIST then
      flush_run()
      items[#items + 1] = { kind = 3, first = n, last = n, w = n.width, rigid = true }
    else
      attach(n)
    end
    n = nxt
  end
  flush_run()
  if lead and items[1] and items[1].first ~= lead then items[1].first = lead end
  return items, indent, strut_h, strut_d
end

-- 自然宽（账本）：每个条目自己的宽 nw，与前一条目之间的隙 gap；顺带标出压缩位权重 spot
local function natural(items, P)
  local p = P.cell
  for i, it in ipairs(items) do
    local prev = items[i - 1]
    local nx = items[i + 1]
    local nk = nx and nx.kind
    local gap = 0
    if it.kind == 0 then
      it.nw = p
    elseif it.kind == 1 then
      it.paired = (nk == 1 or nk == 2)
      it.nw = it.paired and p / 2 or p
      -- 后接闭号的只算 0.2 个位，后接开号或汉字的是整位
      it.spot = (nk == 1) and P.pairw or 1
    elseif it.kind == 2 then
      -- 开号后接开号：前一个半宽
      it.paired = (nk == 2)
      it.nw = it.paired and p / 2 or p
      -- 前面是汉字的开号是整位，前面是标点的只算 0.2 个位
      if not prev then it.spot = 0
      elseif prev.kind == 0 then it.spot = 1
      elseif prev.kind == 1 or prev.kind == 2 then it.spot = P.pairw
      else it.spot = 0 end
      it.afterclose = prev and prev.kind == 1
    elseif it.kind == 3 then
      if it.gn then it.nw = it.gw * P.scale else it.nw = it.w end
      if prev and prev.kind == 0 and is_alnum(it.head) then gap = P.auto end
    else
      local cjk = (prev and prev.kind <= 2) or (nx and nk <= 2)
      if cjk then it.nw = p / 2; it.rigid = true else it.nw = P.space end
    end
    if prev and prev.kind == 3 then
      if (it.kind == 0 or it.kind == 2) and is_alnum(prev.tail) then gap = P.auto
      elseif it.kind == 1 and is_alnum(prev.tail) then gap = P.lp end
    end
    it.gap = gap
  end
end

-- 禁则时行内位各自的压缩上限
local function midcap(it, P)
  if M.compat ~= 15 then return it.paired and P.pairdeep or P.cell / 2 end
  if it.kind == 2 then return it.afterclose and 0 or P.deepcap_open end
  return it.paired and P.pairdeep or P.deepcap
end

-- 贪心断行：返回行表 {first,last,width,over,deep,spots}
local function decide(items, P, measure, cellcap, indent)
  natural(items, P)
  local p = P.cell
  local lines = {}
  local n = #items
  local capacity = measure
  if M.compat ~= 15 and not M.ari0 then capacity = cellcap
  elseif M.snap then capacity = cellcap end
  -- 行末的配对半宽标点在行末是整格（后面没有东西跟它配了）
  local function width_of(first, last, isfirst)
    local w = isfirst and indent or 0
    for k = first, last do
      local it = items[k]
      if not (it.kind == 4 and k == first) then
        w = w + ((k == last and it.paired) and p or it.nw) + (k > first and it.gap or 0)
      end
    end
    return w
  end
  -- 压缩位：first..last-1 里的闭号与开号（last 是被塞的项）；返回位表与权重和
  local function spots_of(first, last)
    local sp, K = {}, 0
    for k = first, last - 1 do
      local it = items[k]
      if (it.kind == 1 or it.kind == 2) and (it.spot or 0) > 0 and not (it.kind == 2 and k == first) then
        sp[#sp + 1] = it; K = K + it.spot
      end
    end
    return sp, K
  end
  -- 模式 15 里西文块的扣减：块的「拟合宽」比排出来的宽多 D，扣 D·g(K)；紧挨着被塞字的块不扣，
  -- 第二个起的块只算一半。g 是量出来的表（k=1..10），D 用 '4' 那种块定标（1.38bp）。
  local G = { [1] = 0.5, [2] = 1.05, [3] = 1.05, [4] = 0.95, [5] = 0.9, [6] = 0.85, [7] = 0.8, [8] = 0.75, [9] = 0.7, [10] = 0.65 }
  local function latin_penalty(first, last, K)
    if M.compat ~= 15 or K < 1 then return 0 end
    -- 被塞的字紧跟在西文后面时整行都不扣（探针与范例都这样）
    local before = items[last - 1]
    if before and (before.kind == 3 or before.kind == 4) then return 0 end
    local g = G[math.min(10, math.max(1, math.floor(K + 0.5)))]
    local pen, nblk = 0, 0
    local k = first
    while k < last do
      local it = items[k]
      if it.kind == 3 and it.gn and not it.rigid then
        -- 一个块：连着的西文串与空格
        local w, n = 0, 0
        local e = k
        while e < last and (items[e].kind == 3 or items[e].kind == 4) do
          local x = items[e]
          if x.kind == 3 and x.gn then w = w + x.gw; n = n + x.gn else w = w + x.nw end
          e = e + 1
        end
        -- 只有两侧都是汉字（两边都有汉西隙）的块才扣；挨着标点的量不出扣减
        local pv = items[k - 1]
        if e < last and pv and pv.kind == 0 and items[e].kind == 0 then
          local q = p / 4
          local F = w + 0.48 * P.bp * math.max(0, n - 1) + 8.1 * P.bp
          local R = math.ceil((w + 2 * q) / q - 0.001) * q
          local D = math.max(0, F - R)
          nblk = nblk + 1
          pen = pen + D * g * (nblk == 1 and 1 or 0.5)
        end
        k = e
      else
        k = k + 1
      end
    end
    return pen
  end
  local function budget(K, adj, first, last)
    if M.compat ~= 15 then
      if not M.ari0 then return 0 end
      return (K >= 1 or adj) and p / 2 or 0
    end
    return p * K / (K + 1) - latin_penalty(first, last, K)
  end
  -- 行末是开号：再给开号自己的左半空白；模式 ≤14 里行内每个位再各给半格，至多两个
  local function open_budget(K, first, last)
    if M.compat ~= 15 then
      if not M.ari0 then return 0 end
      return p / 2 + math.min(K, 2) * p / 2
    end
    return p * K / (K + 1) + p / 2
  end
  -- 禁则用的行内位：只算整位的闭号（配对的半宽闭号与开号不算）
  local function mids_of(sp)
    local m = {}
    for _, it in ipairs(sp) do if it.kind == 1 and it.spot >= 1 then m[#m + 1] = it end end
    return m
  end
  local function deep_budget(sp, first, last)
    local m = mids_of(sp)
    local K = #m
    if M.compat == 15 then
      -- 整个闭号挂出去，行内闭号再共出 p·(K+1)/(K+3)；西文块的扣减量出来比塞汉字时大一倍上下
      if K == 0 then return p end
      return p + p * (K + 1) / (K + 3) - 2 * latin_penalty(first, last, K)
    end
    return (M.compat == 11 and p / 2 or p) + math.min(K, 2) * p / 2
  end
  local i = 1
  while i <= n do
    local isfirst = (#lines == 0)
    local first = i
    local last = i - 1
    local j = i
    local deep = false
    while j <= n do
      local it = items[j]
      local w = width_of(first, j, isfirst)
      if w <= capacity + P.eps then
        last = j; j = j + 1
      else
        local need = w - measure
        local prevopen = (j > first) and items[j - 1].kind == 2
        if it.kind == 1 and (it.spot or 0) > 0 or (it.kind == 3 and ascii_close[it.head]) then
          -- 禁则：闭号不能当行首，先压行内位再悬挂；压不进就连前一项一起下行
          local sp = spots_of(first, j)
          if need <= deep_budget(sp, first, j) + P.eps then
            last = j; j = j + 1; deep = true
            while j <= n and items[j].kind == 1 do last = j; j = j + 1 end
          elseif last > first then
            last = last - 1
          end
          break
        end
        local sp, K = spots_of(first, prevopen and j - 1 or j)
        local T = prevopen and open_budget(K, first, j - 1) or budget(K, false, first, j)
        if need <= T + P.eps then
          -- 塞进来了，接着看下一项是不是闭号（禁则）
          last = j; j = j + 1
        else
          if prevopen then last = last - 1 end
          break
        end
      end
    end
    if last < first then last = first; j = first + 1 end
    -- 行尾的空格不要；行尾不能是开号（禁则拉回没成功时），也不能是西文的左括号
    while last > first and (items[last].kind == 4 or items[last].kind == 2 or (items[last].kind == 3 and ascii_open[items[last].tail])) do last = last - 1 end
    if items[last].paired then items[last].nw = p; items[last].paired = false end
    local w = width_of(first, last, isfirst)
    local over = w - measure
    local sp = nil
    if over > 0 then
      sp = spots_of(first, deep and last or last + 1)
      if deep then sp = mids_of(sp) end
    end
    lines[#lines + 1] = { first = first, last = last, width = w, over = over, deep = deep, spots = sp }
    if M.trace then
      local txt = {}
      for k = first, last do local it = items[k]; txt[#txt+1] = it.c and utf8.char(it.c) or (it.kind == 4 and "_" or "[" .. (it.head and utf8.char(it.head) or "?") .. "]") end
      trace("LINE w=%.2f/%.2f over=%.2f deep=%s %s", w/65536, measure/65536, over/65536, tostring(deep), table.concat(txt))
    end
    -- 行首的空格不跳过：它身上可能挂着链接起点这类 whatsit，留在行里、宽度归零
    i = last + 1
  end
  if #lines > 0 then lines[#lines].islast = true end
  return lines
end

-- 渲染一行：定各处胶与 kern 的宽，装成 hlist
local function render(items, line, P, measure, cellcap, strut_h, strut_d, indent, isfirst)
  local p = P.cell
  local a, b = items[line.first], items[line.last]
  local head, last = a.first, b.last
  local nxt = last.next
  last.next = nil
  if nxt then nxt.prev = nil end
  head.prev = nil
  -- 行首的空格不占位
  if a.kind == 4 then a.first.width = 0 end
  -- 所有胶先定死；每个条目的节点跨度补 kern 到账本宽 nw，条目前的隙用 kern 摆
  for g in traverse(head) do
    if g.id == GLUE then g.stretch = 0; g.stretch_order = 0; g.shrink = 0 end
  end
  local function nodewidth(g)
    local id = g.id
    if id == GLYPH or id == HLIST or id == VLIST or id == RULE then return g.width
    elseif id == GLUE then return g.width
    elseif id == KERN then return g.kern
    elseif id == DISC then
      local w = 0
      for h in traverse(g.replace) do if h.id == GLYPH then w = w + h.width elseif h.id == KERN then w = w + h.kern end end
      return w
    else return 0 end
  end
  local autos, spaces, cjkgaps = {}, {}, {}
  for k = line.first, line.last do
    local it = items[k]
    if k > line.first and it.gap ~= 0 then
      local kn = node.new(KERN); kn.kern = it.gap
      node.insert_before(head, it.first, kn)
      if it.gap == P.auto then autos[#autos + 1] = kn end
    end
    if it.kind == 4 then
      it.first.width = (k == line.first) and 0 or it.nw
      if not it.rigid and k > line.first then spaces[#spaces + 1] = it.first end
      -- 空格后面挂着的胶清零
      local g = it.first.next
      while g and g ~= it.last.next do if g.id == GLUE then g.width = 0 end; g = g.next end
    elseif not it.rigid then
      if it.kind == 3 and it.gn then
        -- 西文字宽平衡：串里原有的字距清零，每个字形后补 (scale-1)·宽
        for _, kn in ipairs(it.kerns) do kn.kern = 0 end
        for _, g in ipairs(it.glyphs) do
          local kn = node.new(KERN); kn.kern = g.width * (P.scale - 1)
          node.insert_after(head, g, kn)
          if g == it.last then it.last = kn end
        end
        for _, d in ipairs(it.discs or {}) do
          local kn = node.new(KERN); kn.kern = nodewidth(d) * (P.scale - 1)
          node.insert_after(head, d, kn)
          if d == it.last then it.last = kn end
        end
      end
      local w = 0
      local g = it.first
      while g do w = w + nodewidth(g); if g == it.last then break end; g = g.next end
      local kn = node.new(KERN); kn.kern = it.nw - w
      if it.kind == 2 then
        -- 开号的空白在左边
        node.insert_before(head, it.first, kn)
        if it.first == head then head = kn end
      else
        node.insert_after(head, it.last, kn)
      end
      it.adj = kn
      -- 汉字缝：每个中文条目后面那个补位 kern（行末那个不算，它后面没有字了）
      if it.kind <= 2 and k < line.last then cjkgaps[#cjkgaps + 1] = kn end
    end
  end
  -- 压缩
  local function take(it, amt)
    if amt <= 0 then return end
    if it.kind == 1 then
      it.adj.kern = it.adj.kern - amt
    elseif it.kind == 2 then
      it.adj.kern = it.adj.kern - amt
    end
  end
  if line.over > 0 and line.spots then
    local sp = line.spots
    if not line.deep then
      local K = 0
      for _, it in ipairs(sp) do K = K + it.spot end
      if K > 0 then for _, it in ipairs(sp) do take(it, line.over * it.spot / K) end end
    else
      local rest = line.over
      -- 行内位吃得下：全由行内位按上限比例分担，不挂；
      -- 吃不下：模式 11 挂多少算多少（至多半格），其余模式挂半格或整格（量出来是这两档），
      -- 挂多了行反而松，余量匀给汉字缝；行内位分担挂剩下的那部分。
      -- 模式 ≤14 行内位至多两个，取靠后的两个
      if M.compat ~= 15 and #sp > 2 then
        local m = {}
        for i2 = #sp - 1, #sp do m[#m + 1] = sp[i2] end
        sp = m
      end
      local cap = 0
      for _, it in ipairs(sp) do cap = cap + midcap(it, P) end
      if rest <= cap + P.eps then
        for _, it in ipairs(sp) do take(it, rest * midcap(it, P) / cap) end
        rest = 0
      else
        local r = rest - cap
        local hang
        if M.compat == 11 then hang = math.min(r, p / 2)
        elseif r <= p / 2 + P.eps then hang = p / 2
        else hang = p end
        if hang >= rest then
          -- 挂出去的比缺的还多，行松了
          local extra = hang - rest
          if extra > 0 and #cjkgaps > 0 then
            local each = extra / #cjkgaps
            for _, g in ipairs(cjkgaps) do g.kern = g.kern + each end
          end
          rest = hang
        else
          local r2 = rest - hang
          if cap > 0 then
            local f = math.min(1, r2 / cap)
            for _, it in ipairs(sp) do take(it, f * midcap(it, P)) end
            r2 = r2 - f * cap
          end
          rest = hang + r2
        end
      end
      if rest > 0 then
        local kn = node.new(KERN); kn.kern = -rest
        node.insert_after(head, node.tail(head), kn)
      end
    end
  end
  -- 行尾的胶与罚分去掉
  local t = node.tail(head)
  while t and t ~= head and (t.id == GLUE or t.id == PENALTY) do
    local pr = t.prev; node.remove(head, t); node.free(t); t = pr
  end
  local target = measure
  if line.islast then
    local fil = node.new(GLUE); fil.subtype = 15; fil.stretch = 65536; fil.stretch_order = 2
    node.insert_after(head, node.tail(head), fil)
  elseif line.over <= 0 then
    -- 松行：贴格段伸到整数格，否则伸到版心；余量先给汉西隙与空格，再匀给汉字缝
    if M.snap or (M.compat ~= 15 and not M.ari0) then target = cellcap end
    local slack = target - line.width
    if slack > 0 then
      local n1 = #autos + #spaces
      if n1 > 0 then
        -- 有汉字缝可分时汉西隙与空格每处至多伸 autostretch，纯西文行由空格全包
        local each = slack / n1
        if #cjkgaps > 0 then each = math.min(each, P.autostretch) end
        for _, g in ipairs(autos) do g.kern = g.kern + each end
        for _, g in ipairs(spaces) do g.width = g.width + each end
        slack = slack - each * n1
      end
      if slack > 0 and #cjkgaps > 0 then
        local each = slack / #cjkgaps
        for _, g in ipairs(cjkgaps) do g.kern = g.kern + each end
        slack = 0
      end
    end
    if target < measure then
      local kn = node.new(KERN); kn.kern = measure - target
      node.insert_after(head, node.tail(head), kn)
    end
  end
  if isfirst and indent > 0 then
    local ib = node.new(HLIST); ib.subtype = 3; ib.width = indent
    ib.next = head; head.prev = ib; head = ib
  end
  local box = node.hpack(head, measure, "exactly")
  box.height = strut_h; box.depth = strut_d
  -- 脚注、浮动体、标记、vadjust 不能留在行盒里：照 TeX 断行的做法搬到行盒后面
  local moved = nil
  local g = box.list
  while g do
    local nx = g.next
    if g.id == INS or g.id == MARK or g.id == ADJUST then
      box.list = node.remove(box.list, g)
      g.prev = nil; g.next = nil
      if g.id == ADJUST then
        -- vadjust 的内容是一段竖排材料，整段接上
        local inner = g.list; g.list = nil
        node.free(g)
        if inner then moved = moved and node.insert_after(moved, node.tail(moved), inner) or inner end
      else
        moved = moved and node.insert_after(moved, node.tail(moved), g) or g
      end
    end
    g = nx
  end
  return box, moved
end

function M.filter(head, is_display)
  if (tex.parshape or 0) ~= 0 or tex.hangindent ~= 0 then return true end
  local ls, rs = tex.getskip("leftskip"), tex.getskip("rightskip")
  if ls.width ~= 0 or rs.width > 65536 * 6 or rs.stretch_order > 0 then return true end
  if tex.hsize ~= tex.getdimen("textwidth") then return true end
  if M.attr and (node.get_attribute(head, M.attr) or 0) == 0 then return true end
  local ok, cell = pcall(tex.getdimen, "g_hit_msword_char_dim")
  if not ok or not cell or cell == 0 then return true end
  local em = tex.sp("1em")
  -- 格宽跟着本级字号缩放：正文一格是 g_hit_msword_char_dim，别的字号按比例
  local ok2, body = pcall(tex.getdimen, "g_hit_body_size_dim")
  if ok2 and body and body > 0 then cell = cell * em / body end
  local items, indent, strut_h, strut_d = tokenize(head)
  if #items == 0 then return true end
  local bp = 65781.76
  local P = { cell = cell, em = em, bp = bp, scale = cell / em,
              auto = 3.25 * bp, lp = 0.27 * bp, space = 3.25 * bp,
              pairw = 0.2, deepcap = cell / 2, deepcap_open = 5.5 * bp, pairdeep = 2.0 * bp,
              hang = cell / 2, autostretch = 2.8 * bp,
              eps = 0.03 * bp }
  local measure = tex.hsize
  local cellcap = math.floor(measure / cell) * cell
  local lines = decide(items, P, measure, cellcap, indent)
  local newhead, tail = nil, nil
  local prevdepth = tex.prevdepth
  for li, line in ipairs(lines) do
    local box, moved = render(items, line, P, measure, cellcap, strut_h, strut_d, indent, li == 1)
    if li > 1 then
      local g = node.new(GLUE); g.subtype = 2
      g.width = tex.baselineskip.width - prevdepth - box.height
      if g.width < 0 then g.width = 0 end
      tail.next = g; g.prev = tail; tail = g
    end
    if newhead == nil then newhead = box else tail.next = box; box.prev = tail end
    tail = box
    if moved then
      tail.next = moved; moved.prev = tail; tail = node.tail(moved)
    end
    prevdepth = box.depth
  end
  return newhead
end

-- 装回调。attr 是类分的属性号：段首节点上为 0 的段不管（交回 TeX）
function M.install(attr)
  M.attr = attr
  luatexbase.add_to_callback("linebreak_filter",
    function(h, d) return M.filter(h, d) end, "hit.linebreak")
end
function M.uninstall()
  if luatexbase.in_callback("linebreak_filter", "hit.linebreak") then
    luatexbase.remove_from_callback("linebreak_filter", "hit.linebreak")
  end
end

return M
%</lua>
%    \end{macrocode}
%
% \Finale
